\paragraph{bin-fold escort 温度流形的一维 Bernoulli 塌缩、有限弧长相图与全阶累积量闭包}
在末位充分化、两点 \texorpdfstring{$f$}{f}-散度闭式与 Fisher--Rao 距离都已确定之后，bin-fold 的 escort 温度族还可进一步被压缩为一条真正的一维极限统计流形：其尾分位阈值与高阶涨落都由同一 Bernoulli 末位律完全锁定。

\begin{theorem}[末位比特把 escort 温度族识别为一维 Bernoulli \texorpdfstring{$f$}{f}-散度流形]\label{thm:xi-time-part9ob-escort-lastbit-bernoulli-fdiv}
沿用结论 \ref{con:xi-time-part9o-escort-fdiv-collapse} 与推论 \ref{cor:xi-time-part9l-escort-two-block-normal-form} 的记号，并记
\[
\ell(w):=w_m,
\qquad
\theta(q):=\frac{1}{1+\varphi^{q+1}}.
\]
则对任意 \(p,q\ge 0\) 与任意凸函数 \(f:(0,\infty)\to \mathbb R\)（满足 \(f(1)=0\)），都有
\[
\lim_{m\to\infty}D_f(\pi_{m,q}\Vert \pi_{m,p})
=
\bigl(1-\theta(p)\bigr)\,
f\!\Bigl(\frac{1-\theta(q)}{1-\theta(p)}\Bigr)
\;+\;
\theta(p)\,
f\!\Bigl(\frac{\theta(q)}{\theta(p)}\Bigr).
\]
换言之，整个 escort 温度族在全部 Csisz\'ar--Morimoto \(f\)-散度意义下都塌缩到一维 Bernoulli 统计流形
\[
q\longmapsto \mathrm{Bernoulli}(\theta(q)),
\]
而末位比特 \(\ell\) 正是该族的极限充分统计量。
\end{theorem}
\begin{proof}
上述 \(f\)-散度闭式正是结论 \ref{con:xi-time-part9o-escort-fdiv-collapse} 的内容。

另一方面，推论 \ref{cor:xi-time-part9l-escort-two-block-normal-form} 给出
\[
\pi_{m,q}
=
\bigl(1-\theta(q)\bigr)u_{A_0}
\;+\;
\theta(q)u_{A_1}
\;+\;
\varepsilon_{m,q},
\]
其中
\[
A_b:=\{w\in X_m:\ w_m=b\},
\qquad
u_{A_b}:=U_{m,b},
\qquad
\|\varepsilon_{m,q}\|_{\mathrm{TV}}
=
O\!\left(\Bigl(\frac{\varphi}{2}\Bigr)^m\right).
\]
因此，对每个固定 \(q\ge 0\)，\(\pi_{m,q}\) 以指数小误差退化为由末位读出 \(\ell\) 所控制的两块均匀混合，而该混合正是 \(\mathrm{Bernoulli}(\theta(q))\) 经提升核 \(L_m(b,\cdot)=u_{A_b}\) 的输出。故末位比特 \(\ell\) 在极限中成为充分统计量；再与前述 \(f\)-散度闭式合并，即得所述一维 Bernoulli 统计流形解释。证毕。
\end{proof}

\begin{theorem}[escort 温度流形的 Fisher--Rao 完整几何]\label{thm:xi-time-part9ob-escort-fisher-total-length}
沿用结论 \ref{con:xi-time-part9o-escort-fisher-rao-closed} 的记号。定义极限 Fisher 信息密度
\[
\mathcal I_\infty(q)
:=
\lim_{p\to q}
\frac{2\,D_{\mathrm{KL}}(\mathrm{Bernoulli}(\theta(q))\Vert \mathrm{Bernoulli}(\theta(p)))}{(q-p)^2}.
\]
则
\[
\boxed{\;
\mathcal I_\infty(q)
=
(\log\varphi)^2\,\theta(q)\bigl(1-\theta(q)\bigr)
=
(\log\varphi)^2\frac{\varphi^{q+1}}{(1+\varphi^{q+1})^2}.\;}
\]
相应的 Fisher--Rao 度量可写成
\[
\mathrm ds^2=\mathcal I_\infty(q)\,\mathrm dq^2,
\]
\[
d_{\mathrm{FR}}(q_1,q_2)
=
2\left|
\arcsin\sqrt{\theta(q_2)}-\arcsin\sqrt{\theta(q_1)}
\right|.
\]
进一步，整条温度射线 \(q\in[0,\infty)\) 的总信息长度有限，并且
\[
\mathcal L_{[0,\infty)}
:=
\int_0^\infty \sqrt{\mathcal I_\infty(q)}\,\mathrm dq
=
2\arctan(\varphi^{-1/2})
=
2\arcsin(\varphi^{-1}).
\]
\leanpartial{paper\_xi\_time\_part9ob\_escort\_fisher\_total\_length}{The requested Lean signature quantifies over arbitrary functions and constants, so \texttt{paper\_xi\_time\_part9ob\_escort\_fisher\_total\_length} is not derivable without additional hypotheses.}
\end{theorem}
\begin{proof}
由结论 \ref{con:xi-time-part9o-escort-fisher-rao-closed}，Bernoulli 极限模型的 Fisher--Rao 度量正由
\[
\frac{(\theta'(q))^2}{\theta(q)(1-\theta(q))}
\]
给出，因此
\[
\mathcal I_\infty(q)
=
\frac{(\theta'(q))^2}{\theta(q)(1-\theta(q))}.
\]
再由
\[
\theta(q)=\frac{1}{1+\varphi^{q+1}}
\]
直接求导可得
\[
\theta'(q)
=
-(\log\varphi)\,\theta(q)\bigl(1-\theta(q)\bigr),
\]
故
\[
\mathcal I_\infty(q)
=(\log\varphi)^2\,\theta(q)\bigl(1-\theta(q)\bigr)
=(\log\varphi)^2\frac{\varphi^{q+1}}{(1+\varphi^{q+1})^2}.
\]
这就给出 Fisher 信息密度闭式。测地距离公式则与结论 \ref{con:xi-time-part9o-escort-fisher-rao-closed} 完全一致。

对总长度，令
\[
a:=\varphi^{q+1},
\qquad
\mathrm dq=\frac{\mathrm da}{a\log\varphi}.
\]
于是
\[
\sqrt{\mathcal I_\infty(q)}
=(\log\varphi)\frac{\sqrt a}{1+a},
\]
从而
\[
\mathcal L_{[0,\infty)}
=
\int_0^\infty \sqrt{\mathcal I_\infty(q)}\,\mathrm dq
=
\int_{\varphi}^{\infty}\frac{a^{-1/2}}{1+a}\,\mathrm da.
\]
再令 \(a=t^2\)，则 \(\mathrm da=2t\,\mathrm dt\)，故
\[
\mathcal L_{[0,\infty)}
=
2\int_{\sqrt\varphi}^{\infty}\frac{1}{1+t^2}\,\mathrm dt
=
\pi-2\arctan(\sqrt\varphi).
\]
由恒等式
\[
\arctan x+\arctan(x^{-1})=\frac{\pi}{2}
\qquad (x>0)
\]
可得
\[
\pi-2\arctan(\sqrt\varphi)=2\arctan(\varphi^{-1/2}).
\]
另一方面，由
\[
\sin\!\bigl(2\arctan(\varphi^{-1/2})\bigr)
=
\frac{2\varphi^{-1/2}}{1+\varphi^{-1}}
=
\varphi^{-1},
\]
且 \(2\arctan(\varphi^{-1/2})\in(0,\pi/2)\)，故
\[
2\arctan(\varphi^{-1/2})=2\arcsin(\varphi^{-1}).
\]
证毕。
\end{proof}

\begin{theorem}[escort 尾分位阈值的完整相图由显式相变曲线控制]\label{thm:xi-time-part9ob-escort-quantile-phase-diagram}
对任意 \(q\ge 0\)，令
\[
W_m^{(q)}\sim \pi_{m,q},
\]
并定义尾部分位指数
\[
\gamma_{m,q}^{\mathrm{bin}}(\varepsilon)
:=
\inf\left\{
\gamma:\ 
\mathbf P_{\pi_{m,q}}\!\left(
\frac{1}{m}\log d_m^{\mathrm{bin}}(W_m^{(q)})>\gamma
\right)\le \varepsilon
\right\},
\qquad
0<\varepsilon<1,
\]
以及相应阈值
\[
B_{m,q}^{\mathrm{bin}}(\varepsilon)
:=
\exp\!\bigl(m\gamma_{m,q}^{\mathrm{bin}}(\varepsilon)\bigr).
\]
再定义显式相变曲线
\[
\varepsilon_{\mathrm c}(q)
:=
1-\theta(q)
=
\frac{\varphi^{q+1}}{1+\varphi^{q+1}}.
\]
则对任意固定 \(q\ge 0\) 与任意
\[
\varepsilon\in(0,1)\setminus\{\varepsilon_{\mathrm c}(q)\},
\]
都有
\[
\boxed{\;
\left(\frac{\varphi}{2}\right)^m B_{m,q}^{\mathrm{bin}}(\varepsilon)
\longrightarrow
\begin{cases}
1, & 0<\varepsilon<\varepsilon_{\mathrm c}(q),\\[4pt]
\varphi^{-1}, & \varepsilon_{\mathrm c}(q)<\varepsilon<1.
\end{cases}
\;}
\]
此外，相变曲线完全显式，且严格递增：
\[
\varepsilon_{\mathrm c}'(q)
=
\frac{\log\varphi\ \varphi^{q+1}}{(1+\varphi^{q+1})^2}
>
0.
\]
当 \(q=1\) 时，上式退化为推论 \ref{cor:fold-bin-quantile-threshold-constant} 的原始 bin-fold 阈值常数律。
\leanpartial{paper\_xi\_time\_part9ob\_escort\_quantile\_phase\_diagram}{The requested Lean signature quantifies over arbitrary functions and constants, so \texttt{paper\_xi\_time\_part9ob\_escort\_quantile\_phase\_diagram} is not derivable without additional hypotheses.}
\end{theorem}
\begin{proof}
定义缩放随机变量
\[
Z_m^{(q)}
:=
\left(\frac{\varphi}{2}\right)^m d_m^{\mathrm{bin}}(W_m^{(q)}).
\]
由定理 \ref{thm:fold-bin-two-state-asymptotic} 的一致展开，
\[
Z_m^{(q)}
=
\varphi^{-W_m^{(q)}(m)}
\left(1+O\!\left(\Bigl(\frac{\varphi}{2}\Bigr)^m\right)\right),
\]
其中误差在 \(X_m\) 上一致。另一方面，由命题 \ref{prop:fold-bin-escort-lastbit}，
\[
\mathbf P_{\pi_{m,q}}\bigl(W_m^{(q)}(m)=1\bigr)
=
\theta(q)+O\!\left(\Bigl(\frac{\varphi}{2}\Bigr)^m\right),
\]
\[
\mathbf P_{\pi_{m,q}}\bigl(W_m^{(q)}(m)=0\bigr)
=
1-\theta(q)+O\!\left(\Bigl(\frac{\varphi}{2}\Bigr)^m\right).
\]
故
\[
Z_m^{(q)}\Longrightarrow Z^{(q)},
\]
其中
\[
\mathbf P\bigl(Z^{(q)}=1\bigr)=1-\theta(q)=\varepsilon_{\mathrm c}(q),
\qquad
\mathbf P\bigl(Z^{(q)}=\varphi^{-1}\bigr)=\theta(q).
\]

而
\[
\left(\frac{\varphi}{2}\right)^m B_{m,q}^{\mathrm{bin}}(\varepsilon)
\]
正是随机变量 \(Z_m^{(q)}\) 在同一尾部分位定义下的阈值。由于
\[
\varphi^{-1}<1,
\]
极限两点分布的上尾量在 \(\varepsilon=\varepsilon_{\mathrm c}(q)\) 处发生唯一跳变：若 \(\varepsilon<\varepsilon_{\mathrm c}(q)\)，则必须把原子 \(1\) 一并压掉，故极限阈值为 \(1\)；若 \(\varepsilon>\varepsilon_{\mathrm c}(q)\)，则容许丢弃原子 \(1\) 的全部质量，故极限阈值降为 \(\varphi^{-1}\)。这便得到所示分段极限。

最后，
\[
\varepsilon_{\mathrm c}(q)=\frac{\varphi^{q+1}}{1+\varphi^{q+1}}
\]
直接求导即得
\[
\varepsilon_{\mathrm c}'(q)
=
\frac{\log\varphi\ \varphi^{q+1}}{(1+\varphi^{q+1})^2}>0.
\]
当 \(q=1\) 时，
\[
\varepsilon_{\mathrm c}(1)
=
\frac{\varphi^2}{1+\varphi^2}
=
\frac{\varphi}{\sqrt5},
\]
故恰好恢复推论 \ref{cor:fold-bin-quantile-threshold-constant}。证毕。
\end{proof}

\begin{theorem}[escort 纤维信息极限的全阶累积量闭式与永久负偏斜]\label{thm:xi-time-part9ob-escort-loginfo-cumulants-negative-skew}
对任意 \(q\ge 0\)，令 \(W_m^{(q)}\sim \pi_{m,q}\)，并定义中心化纤维信息
\[
Y_m^{(q)}
:=
\log d_m^{\mathrm{bin}}(W_m^{(q)})
-m\log\!\Bigl(\frac{2}{\varphi}\Bigr).
\]
则
\[
Y_m^{(q)}\Longrightarrow Y^{(q)},
\]
其中极限分布是两点律
\[
Y^{(q)}=
\begin{cases}
0,& \text{概率 }1-\theta(q),\\
-\log\varphi,& \text{概率 }\theta(q).
\end{cases}
\]
因此其累积母函数为
\[
K_q(t)
:=
\log \mathbf E\!\left[e^{tY^{(q)}}\right]
=
\log\!\bigl(1-\theta(q)+\theta(q)\varphi^{-t}\bigr),
\qquad t\in\mathbb R.
\]
若记
\[
\kappa_r^{(q)}:=K_q^{(r)}(0),
\]
则前四阶累积量满足
\[
\kappa_1^{(q)}=-\theta(q)\log\varphi,
\]
\[
\kappa_2^{(q)}=(\log\varphi)^2\,\theta(q)\bigl(1-\theta(q)\bigr),
\]
\[
\kappa_3^{(q)}=-(\log\varphi)^3\,\theta(q)\bigl(1-\theta(q)\bigr)\bigl(1-2\theta(q)\bigr),
\]
\[
\kappa_4^{(q)}=(\log\varphi)^4\,\theta(q)\bigl(1-\theta(q)\bigr)\bigl(1-6\theta(q)+6\theta(q)^2\bigr).
\]
特别地，对所有 \(q\ge 0\) 都有
\[
\kappa_3^{(q)}<0.
\]
当 \(q=1\) 时，上述闭式恰化为定理 \ref{thm:xi-fold-hidden-logmultiplicity-laplace-cumulant-closed} 的均匀基线公式。
\end{theorem}
\begin{proof}
由定理 \ref{thm:fold-bin-two-state-asymptotic} 的一致对数展开，
\[
\log d_m^{\mathrm{bin}}(w)
=
m\log\!\Bigl(\frac{2}{\varphi}\Bigr)-w_m\log\varphi
+O\!\left(\Bigl(\frac{\varphi}{2}\Bigr)^m\right)
\]
在 \(w\in X_m\) 上一致成立，因此
\[
Y_m^{(q)}
=
-W_m^{(q)}(m)\log\varphi
+O_{\mathbf P}\!\left(\Bigl(\frac{\varphi}{2}\Bigr)^m\right).
\]
再与命题 \ref{prop:fold-bin-escort-lastbit} 的末位极限
\[
\mathbf P_{\pi_{m,q}}\bigl(W_m^{(q)}(m)=1\bigr)\to \theta(q)
\]
合并，即得所示两点极限分布。

于是对任意实数 \(t\)，有
\[
\mathbf E\!\left[e^{tY^{(q)}}\right]
=
\bigl(1-\theta(q)\bigr)+\theta(q)\varphi^{-t},
\]
取对数便得到 \(K_q(t)\)。

其后各阶累积量只需在 \(t=0\) 处逐次求导；等价地，也可写成
\[
Y^{(q)}=-(\log\varphi)\,B_q,
\qquad
B_q\sim \mathrm{Bernoulli}(\theta(q)),
\]
并直接调用 Bernoulli 两点律的标准累积量公式，即得所列四式。

最后，由
\[
0<\theta(q)\le \theta(0)=\frac{1}{1+\varphi}=\varphi^{-2}<\frac12
\]
可知
\[
1-2\theta(q)>0,
\]
故
\[
\kappa_3^{(q)}
=
-(\log\varphi)^3\,\theta(q)\bigl(1-\theta(q)\bigr)\bigl(1-2\theta(q)\bigr)<0.
\]
这便给出永久负偏斜性。证毕。
\end{proof}

\begin{theorem}[信息谱的二层量子化]\label{thm:xi-time-part9ob-escort-surprisal-two-level-quantization}
对任意 \(q\ge 0\)，令
\[
W_m^{(q)}\sim \pi_{m,q},
\qquad
\ell_m(w):=w_m,
\qquad
I_m^{(q)}:=-\log \pi_{m,q}(W_m^{(q)}).
\]
并记
\[
c(q):=\log\!\left(\frac{\varphi+\varphi^{-q}}{\sqrt5}\right).
\]
则有概率展开
\[
I_m^{(q)}
=
m\log\varphi+c(q)+q\,\ell_m(W_m^{(q)})\log\varphi+o_{\mathbf P}(1).
\]
因此
\[
I_m^{(q)}-m\log\varphi-c(q)
\Longrightarrow
q\log\varphi\cdot B_q,
\qquad
B_q\sim \mathrm{Bernoulli}(\theta(q)).
\]
特别地，在原始分布 \(\mu_m=\pi_{m,1}\) 下，
\[
I_m^{(1)}-m\log\varphi
\Longrightarrow
\log\varphi\cdot B,
\]
\[
\mathbf P(B=1)=\frac{1}{\varphi\sqrt5},
\qquad
\mathbf P(B=0)=\frac{\varphi}{\sqrt5}.
\]
\end{theorem}
\begin{proof}
由定理 \ref{thm:fold-bin-two-state-asymptotic} 与命题 \ref{prop:fold-bin-power-sum-asymptotic}，对 \(w\in X_m\) 一致有
\[
\pi_{m,q}(w)
=
\frac{\varphi^{-q\ell_m(w)}}{F_{m+1}+\varphi^{-q}F_m}
\Bigl(1+O\!\left(\Bigl(\frac{\varphi}{2}\Bigr)^m\right)\Bigr).
\]
取负对数得
\[
-\log \pi_{m,q}(w)
=
\log(F_{m+1}+\varphi^{-q}F_m)
+q\,\ell_m(w)\log\varphi
+O\!\left(\Bigl(\frac{\varphi}{2}\Bigr)^m\right),
\]
且误差在 \(w\) 上一致。另一方面，由 Binet 公式，
\[
F_{m+1}+\varphi^{-q}F_m
=
\frac{\varphi^m}{\sqrt5}\bigl(\varphi+\varphi^{-q}\bigr)+O(\varphi^{-m}),
\]
从而
\[
\log(F_{m+1}+\varphi^{-q}F_m)
=
m\log\varphi+c(q)+o(1).
\]
将两式合并并代入 \(w=W_m^{(q)}\)，即得
\[
I_m^{(q)}
=
m\log\varphi+c(q)+q\,\ell_m(W_m^{(q)})\log\varphi+o_{\mathbf P}(1).
\]
再由命题 \ref{prop:fold-bin-escort-lastbit}，
\[
\ell_m(W_m^{(q)})\Longrightarrow \mathrm{Bernoulli}(\theta(q)),
\]
故第二个极限立即成立。

当 \(q=1\) 时，
\[
c(1)=\log\!\left(\frac{\varphi+\varphi^{-1}}{\sqrt5}\right)=0,
\]
且
\[
\theta(1)=\frac{1}{1+\varphi^2}=\frac{1}{\varphi\sqrt5},
\qquad
1-\theta(1)=\frac{\varphi}{\sqrt5}.
\]
于是得到最后两式。证毕。
\end{proof}

\endinput
